\section{Science and Wisdom}

Despite is lofty but overrated reputation in our day, science is actually a rather low form of knowledge. First of all, it is a product of the rational, discursive mind, the ``soul'' or ``mind'', whereas authentic spiritual understanding comes from the intellect, the nous, the ``spirit''.

This is exemplified by the democratic nature of scientific phenomena. It is assumed that everyone with normal sensory apparatus experiences the world more or less the same, no matter what his state of development. Science assumes this universality and through various external experiments creates conceptual schemes assumed to account to the phenomena. However, its abductive logic can only produce what Plato called ``a likely story'', always open to change and development. In more recent times, this has been stated most forcibly by Karl Popper who logically demonstrated that a scientific theory can never be infallibly established.

Yet, today many fall under the illusion that Quantum Theory somehow reveals some spiritual truths. Yet the only useful aspect of Quantum Theory is quantitative, which yields quite accurate results in the observation of certain phenomena. There are at least a half dozen conceptual schemes to provide a qualitative description of quantum phenomena with no scientific way to select from among them. To choose a more or less ``spiritual'' explanation is purely arbitrary and even fanciful, certainly not ``scientific''. If any truly spiritual insights have come forth from the many popularisers of Quantum Theory, then their true source is actually elsewhere and in no way are attributable to ``science''.

Its foundation rests on Schroedinger's equation which reduces all ``matter'' to probability waves in a state of inderminateness until it encounters the presence of a ``conscious'' observer. Since the equation is linear, the entire material/energy universe can be represented by such an equation. Once that step has been take, it then becomes impossible to reconstruct the world of objective things from physics alone. True, this is more or less what would be expected from the Traditional point of view, that Prime Matter is chaos and without form, until informed by consciousness. However, that is no longer physics, but metaphysics.

Authentic spiritual knowledge, on the other hand, arises in the spirit of man, and is given the name ``Wisdom''. Science observes phenomena and posits some other realty behind it. Wisdom in contrast is direct and intuitive, hence certain. Science in principle is open to all with the proper education --- no other self-transformation is required.

Wisdom is for the few who make the required effort. Wisdom requires a change in one's state of being, that is, there is no distinction between knowledge and the thing known. In other words, to understand higher states, it is necessary to be transformed into that higher state.

Of course, this goes against the grain of all modernity by offending its democratic pretensions. The modern man wants to read a book, attend a lecture, take notes, and then become wise.

Few are those who are willing to travel the road to Wisdom, and fewer still are those who can show the way.


\flrightit{Posted on 2007-11-19 by Cologero}

